\section{Introduction}
The quantum Fisher information matrix (QFIM) equips a parameterized quantum state with a local information geometry and underlies the quantum natural-gradient update \cite{Braunstein1994,Stokes2020}. This geometry is measurement independent: it quantifies local distinguishability available in the state before a particular experimental interface is fixed. Variational quantum learning, by contrast, is normally driven by a restricted classical record, for example Pauli expectation values, low-order correlators, or marginals. Consequently, the optimizer need not have operational access to every direction represented by the full QFIM.

Our earlier work made this separation explicit. Under joint state--tangent isotropy, a fixed basis followed by a rank-$r$ diagonal readout retains a dimension-controlled fraction of the visible score information \cite{AitHaddouRank}. The subsequent spectral theory removed the isotropy requirement and represented measurement/readout accessibility by a positive contraction $T_{\mathcal M,\mathcal A}$ acting on normalized quantum tangent space \cite{AitHaddouSpectral}. If $C_Q$ denotes the orientation covariance of normalized quantum tangents, the ensemble-averaged accessible fraction is
\begin{equation}
 \mathbb E\!\left[\frac{I_{\mathcal M,\mathcal A}}{F_Q}\right]=\Tr(T_{\mathcal M,\mathcal A}C_Q).
 \label{eq:master}
\end{equation}
The isotropic readout-rank law is recovered as a special case when $C_Q$ is proportional to the identity on the relevant tangent space.

Equation~\eqref{eq:master} suggests an optimization question distinct from estimating the full QFIM: if the learning interface exposes only a contracted tangent geometry, should the optimizer precondition the loss gradient with directions that the interface does not retain? Here we answer this question constructively by introducing the \emph{accessible quantum natural gradient} (AQNG). AQNG is not defined as a numerical approximation whose sole objective is to reproduce the full QNG matrix. Instead, it uses the metric induced by retained measurement/readout tangent features. An accessible metric may carry only a modest fraction of QFIM trace and nevertheless generate a useful optimization direction.

We test this principle on four binary classification tasks and four circuit architectures. AQNG and full QNG share the same initial circuit, data split, stochastic seed, loss gradient, and training budget within each comparison. The results reveal three regimes. SU(2)-Haar provides a positive case in which a low-trace accessible metric is sufficient and frequently superior in generalization. SU(2)-CNOT and RyRz-CZ show architecture- and dataset-dependent behavior rather than universal dominance. Finally, a number-conserving $U(1)$ circuit provides a negative control: AQNG and full-QNG update directions are highly aligned, yet both remain at fixed chance-like classification accuracy. Natural-gradient fidelity therefore cannot compensate for a task-misaligned model/readout interface.
